\documentclass[UTF8,fontset=none,11pt,a4paper]{ctexart}
\usepackage{ipara-handbook}
\handbooksetup
  {DS-03 串与模式匹配}
  {408 · 数据结构 DS-03}
  {String · Prefix · Mismatch · Reuse}
\begin{document}
\handbooktitle
  {串与模式匹配}
  {失配之后，怎样复用已经确认的前缀信息}
  {408 · 数据结构 Topic DS-03}

\begin{corebox}{母问题：失配之后，哪些信息真的已经知道？}
给定文本 $T$ 和模式 $P$，匹配不是“把两个字符串碰运气比较”，而是一个不断维护匹配状态的扫描过程：
\[
\begin{aligned}
\text{Matched Prefix}&\to\text{Mismatch}\to\text{Reusable Border}\\
&\to\text{Next Alignment}\to\text{Continue / Accept}.
\end{aligned}
\]
朴素算法在失配后把模式移回下一个起点，因此反复比较文本；KMP 观察已匹配的模式前缀，寻找同时也是后缀的最长边界（border），只缩短模式状态。箭头表示一次匹配状态机的时间顺序，不是某个特定 `next` 数组下标的背诵顺序。
\end{corebox}

\begin{methodbox}{本册定位与 Owner 边界}
\begin{itemize}
\item \kw{Owns：}串的顺序对象、朴素匹配、prefix function/failure 信息、失配回退与 KMP 主循环。
\item \kw{Uses：}Atlas 的表示—操作—不变量—成本语言；底层字符序列只借用一般线性存储，不重新定义 DS01。
\item \kw{Outputs：}向算法专题提供线性时间模式匹配的可调用接口。
\item \kw{Stops：}完整 DFA/自动机最小化、多模式匹配与文本索引属于 Extension 或其他专题。
\end{itemize}
\end{methodbox}

\tableofcontents\newpage

\section{对象、契约与第一基线}
串是一个有序字符序列；模式匹配的输出契约是：返回 $P$ 在 $T$ 中第一次出现的起始位置，若不存在则返回 `npos`。本册代码另约定空模式返回 0，这使边界行为显式而不依赖调用者猜测。

朴素匹配枚举起点 $s=0,1,\ldots,|T|-|P|$，在每个起点从左到右比较。若在第 $j$ 个字符失配，已经比较过的前缀信息只用于判定当前起点失败，算法不会把它带到下一起点。最坏比较次数可以达到 $\Theta((n-m+1)m)$，其中 $n=|T|,m=|P|$。

\begin{examplebox}{为什么需要新机制}
文本 $T=\texttt{aaaaab}$、模式 $P=\texttt{aaab}$ 时，许多起点都会先匹配若干个 `a` 再失配。朴素算法每次都重新确认这些 `a`；但“已经看到的后缀由若干 `a` 构成”并没有消失。若模式自身存在重复前后缀，就可以把这份信息压缩成下一次应比较的模式位置。
\end{examplebox}

\section{Prefix Function：模式自己的可复用结构}
对模式前缀 $P[0..i]$，定义 $\pi[i]$ 为其最长真前缀且同时为后缀的长度。它只由模式自身决定，不依赖文本：
\[
\pi[i]=\max\{k< i+1:\;P[0..k-1]=P[i-k+1..i]\}.
\]
若当前边界长度为 $b=\pi[i-1]$ 且新字符不匹配，任何更短的候选边界都必然是长度 $b$ 边界的边界，因此可令 $b\leftarrow\pi[b-1]$，继续尝试，而不用从 0 重新比较。

这给出构造不变量：处理 $P[0..i]$ 后，`prefix[i]` 恰是该前缀的最长 border；回退链严格变短，所以整个构造是线性的摊还过程。

\subsection{把 \texorpdfstring{$\pi$}{pi} 真正算一遍：\code{ababaca}}
逐个加入新字符，并始终问“当前最长 border 能否延长；不能时，它自己的最长 border 是谁”：
\begin{center}\small
\begin{tabularx}{\linewidth}{C{0.8cm}C{1.1cm}L{2.7cm}Y C{1.1cm}}
\toprule
$i$ & $P[i]$ & 旧候选 $b$ & 比较、回退与可复用事实 & $\pi[i]$\\
\midrule
0 & a & 无 & 单字符没有非空真 border & 0\\
1 & b & 0 & $b\ne a$，无法延长 & 0\\
2 & a & 0 & $a=a$，空 border 延长为 `a' & 1\\
3 & b & 1 & $b=P[1]$，`a' 延长为 `ab' & 2\\
4 & a & 2 & $a=P[2]$，`ab' 延长为 `aba' & 3\\
5 & c & 3 & $c\ne P[3]$：$3\to1\to0$，仍不等 & 0\\
6 & a & 0 & $a=P[0]$，重新得到 `a' & 1\\
\bottomrule
\end{tabularx}
\end{center}
因此 $\pi=[0,0,1,2,3,0,1]$。第 5 项按 $3\to1\to0$ 回退，是因为长度 3 的 border 已失败，剩余候选只能是这个 border 自己的 border；其他长度不可能同时保持前缀和后缀。

\section{KMP 主循环：模式状态回退，主串不倒退}
扫描文本下标 $i$ 时，变量 `matched` 表示当前文本后缀与模式前缀相等的长度。失配时反复执行
\[
matched\leftarrow\pi[matched-1]
\]
直到匹配、归零或接受完整模式；文本下标 $i$ 仍只前进一次。若 `matched==m`，起点为 $i+1-m$。

\begin{examplebox}{一次完整失配：文本不回头，模式状态缩短}
用 $P=\code{ababaca}$ 扫描文本片段 \code{abababaca}。读完 \code{ababa} 后，$matched=5$；下一文本字符是 `b'，而 $P[5]=\code{c}$。已知文本后缀 \code{aba} 等于模式前缀 \code{aba}，所以令
\[
5\longrightarrow\pi[4]=3.
\]
仍用\kw{同一个文本字符 `b'}比较 $P[3]$，匹配后变成 $matched=4$；主串下标没有倒退，也没有重新比较已经证实的 \code{aba}。这就是代码中 `while` 回退后紧接 `if` 比较的语义。
\end{examplebox}

\begin{center}\small
\begin{tabularx}{\linewidth}{L{2.3cm}YY}
\toprule
状态 & 迁移 & 必须保持的事实\\
\midrule
\code{matched=0} & 比较当前文本字符与 $P[0]$ & 没有可复用的模式前缀\\
\code{matched>0} 且失配 & 沿 prefix 链缩短模式状态 & 不移动文本下标，保留已证实的后缀\\
匹配后 & \code{matched++} & 当前状态代表长度为 \code{matched} 的相等前缀\\
\code{matched=m} & 返回 $i+1-m$ & 所有模式字符已连续匹配\\
\bottomrule
\end{tabularx}
\end{center}

\section{代码：把两种扫描策略放在同一契约下}
完整 C++17 实现位于 \codepath{code/ds03_string_matching.hpp}，测试位于 \codepath{code/ds03_string_matching_test.cpp}；正文代码块直接从真实源文件抽取。

\subsection{朴素基线：每个起点重新建立 matched}
\lstinputlisting[
  language=C++,firstline=10,lastline=27,firstnumber=10,numbers=left,
  caption={朴素匹配的起点枚举与局部比较}
]{30_408/10_数据结构/03_串与模式匹配/code/ds03_string_matching.hpp}

朴素版本的重要性在于提供复杂度基线：它没有错误，只是没有利用跨起点的结构信息。任何声称 KMP 更快的解释，都应能指出它省掉了哪类重复比较。

\subsection{Prefix function：失配时只沿 border 链回退}
\lstinputlisting[
  language=C++,firstline=30,lastline=47,firstnumber=30,numbers=left,
  caption={prefix function 的回退链构造}
]{30_408/10_数据结构/03_串与模式匹配/code/ds03_string_matching.hpp}

`border` 只表示当前候选前缀长度。每次回退都跳到更短 border，因此不会形成无限循环；若新字符与候选前缀末字符相等，再增长一次。prefix 数组是 KMP 运行阶段可复用信息的唯一来源。

\subsection{KMP：失配不倒退主串}
\lstinputlisting[
  language=C++,firstline=49,lastline=72,firstnumber=49,numbers=left,
  caption={KMP 主循环与完整模式接受}
]{30_408/10_数据结构/03_串与模式匹配/code/ds03_string_matching.hpp}

主串循环变量 `i` 单调递增；模式状态 `matched` 才可能沿 prefix 链下降。于是预处理 $O(m)$ 加扫描 $O(n)$，总时间 $O(n+m)$，额外空间 $O(m)$。这不是“next 公式自动带来线性”，而是由两个单调性共同保证。

\subsection{测试：重复前后缀与边界返回值}
\lstinputlisting[
  language=C++,firstline=9,lastline=45,firstnumber=9,numbers=left,
  caption={空模式、重复前后缀、失配回退与基线一致性}
]{30_408/10_数据结构/03_串与模式匹配/code/ds03_string_matching_test.cpp}

测试要求朴素与 KMP 在多组文本/模式上返回同一首个位置，直接检查 prefix 数组的重复结构，并覆盖空模式、模式过长、完全无匹配和失配后仍有答案的情形。

测试还枚举二字母表上长度不超过 6 的全部文本、长度不超过 4 的全部模式，逐对验证 KMP 与朴素算法返回同一首个位置。穷举不是正确性证明，但会系统攻击空模式、重叠匹配、周期串、长 border 链与无匹配状态。

\begin{lstlisting}[language=bash]
clang++ -std=c++17 -Wall -Wextra -Werror -pedantic \\
  code/ds03_string_matching_test.cpp -o /tmp/ds03_test
/tmp/ds03_test
\end{lstlisting}

\section{复杂度与边界}
\begin{center}\small
\begin{tabularx}{\linewidth}{L{2.5cm}C{1.8cm}C{1.8cm}YY}
\toprule
任务 & 朴素匹配 & KMP & 关键前提 & 失败边界\\
\midrule
预处理模式 & 无 & $O(m)$ & prefix 语义固定 & 空模式\\
扫描文本 & $O((n-m+1)m)$ 最坏 & $O(n)$ & 主串下标单调前进 & $m>n$\\
总额外空间 & $O(1)$ & $O(m)$ & 只保存模式信息 & 大模式内存预算\\
返回位置 & 首个合法起点 & 首个合法起点 & 接受后立即返回 & 无匹配返回 npos\\
\bottomrule
\end{tabularx}
\end{center}

不同教材的 `next` 可能使用“失配时跳到的下标”或“最长 border 长度”两种编码。做题时先写语义，再把题目下标转换到该语义；不能拿另一版本的初值和边界直接替换。

\begin{methodbox}{从长度语义翻译到教材 next}
本册只维护 $\pi[i]=$“$P[0..i]$ 的最长 border 长度”。若教材把 `next[j]' 定义为“模式状态 $j$ 失配后跳到的状态”，零基长度状态下对应 $\pi[j-1]$（$j>0$）。若教材使用 $-1$ 哨兵或一基下标，先画出“已经匹配多少个字符”，再平移下标；数组首项长什么样不是算法本体。
\end{methodbox}

\section{概念边界与做题控制协议}
\begin{center}\small
\begin{tabularx}{\linewidth}{L{1.9cm}C{0.35cm}L{1.9cm}YY}
\toprule
概念 A & $\neq$ & 概念 B & 真正区别 & 混淆后果\\
\midrule
串 & $\neq$ & 模式 & 被搜索的文本对象 vs 约束模板 & 方向写反\\
border & $\neq$ & 任意重复片段 & 必须同时是前缀和后缀 & 错误回退位置\\
prefix function & $\neq$ & 匹配答案 & 模式内部信息 vs 文本中的出现位置 & 把预处理当搜索\\
失配回退 & $\neq$ & 主串回退 & 只缩短模式状态 & 复杂度退化\\
线性时间 & $\neq$ & 每次比较 $O(1)$ & 总比较次数还需证明单调性 & 复杂度口号化\\
\bottomrule
\end{tabularx}
\end{center}

遇到匹配题，按“固定返回契约 → 写朴素基线 → 找重复前后缀 → 固定 prefix/next 语义 → 画一次失配回退 → 再报告复杂度”的顺序处理。忘掉公式后，只需复原：\kw{已经匹配了什么？失配时哪些后缀仍可能成为模式前缀？模式状态能否缩短而不重扫文本？}

\section{全科坐标与来源处理}
\begin{corebox}{Map Coordinate：Prefix / Failure / Reuse}
DS03 把线性序列上的“匹配失败”转化为模式内部 border 结构；它接收 DS01 的顺序访问成本语言，并向算法专题输出 KMP 的线性扫描接口。自动机、多模式和文本索引是后续 Extension，不改变本册唯一 Owner。
\end{corebox}

本册吸收归档总册对 KMP 的核心表述“已匹配前缀压缩为失败状态”，并将 next 下标差异明确为表示约定问题。代码和测试只验证本册固定的 prefix-length 语义，避免把不同教材的数组编码混作第二套稳定真相。
\end{document}
